% !TeX root = ../main.tex

\chapter{Discusión}
\label{chap:discusion}

Este capítulo interpreta los resultados presentados en el Capítulo~\ref{chap:resultados} y delimita hasta dónde alcanzan. La primera sección examina qué comportamiento describen en conjunto las distintas piezas del análisis, por qué la forma que adopta el fenómeno resulta relevante para la auditoría de un sistema de este tipo y qué queda establecido frente a qué permanece como hipótesis. La segunda revisa los supuestos de los que dependen esas conclusiones y las vías por las que podrían resultar equivocadas. La tercera acota el alcance de lo observado y señala qué extremos requerirían
evidencia adicional. La separación entre amenazas y limitaciones responde a preguntas diferentes: si las inferencias realizadas son correctas dentro del experimento y hasta dónde pueden extenderse fuera de él.

\section{Interpretación de los resultados}
\label{sec:interpretacion_resultados}

El hallazgo central no tiene la forma que cabría esperar. La pregunta habitual ante un sistema de este tipo es si asigna menos capital a las empresas con determinado atributo, y la respuesta es que no hay evidencia de ello: ninguna de las 36 pruebas sobre el \textit{gap} medio sobrevive a la corrección por comparaciones múltiples. Lo que sí aparece es otra cosa. En doce celdas interpretables, la presencia del atributo multiplica por entre tres y ocho la dispersión de las decisiones respecto a la variabilidad basal del sistema. Esto significa que la misma empresa recibe asignaciones muy distintas según la ejecución, y esa inestabilidad depende de si el atributo está presente.

La consecuencia práctica es que una auditoría convencional habría dado este sistema por limpio. Al promediar, las ejecuciones que favorecen a la variante y las que la penalizan se cancelan, y el resultado agregado se parece al de un sistema que trata ambas condiciones por igual. Es la situación que Cooper et al.~\cite{cooper2024arbitrariness} describen al distinguir entre equidad agregada y arbitrariedad de las decisiones, que una métrica media favorable puede convivir con decisiones que no son igualmente reproducibles. El análisis de correlación intraclase sitúa además el fenómeno en la ejecución y no en la cesta, ya que ninguna de las 42 estimaciones sobrevive a la corrección, aunque disponer de solo tres repeticiones por configuración limita la capacidad para detectar si una parte menor del efecto se repite de forma consistente en determinadas cestas.

La composición del comité es el factor que más claramente separa unas condiciones de otras. El exceso de dispersión aparece en tres de las cuatro composiciones homogéneas y en ninguna de las dos heterogéneas. La cuarta homogénea, \textit{homo\_gpt}, tampoco aporta evidencia, aunque por un motivo distinto, porque su referencia ciega resultó inestable y sus celdas quedaron descartadas. El contraste más limpio enfrenta \textit{homo\_llama} con \textit{hetero\_local}, que comparten modelo y rol en el agente ciego y difieren solo en los dos agentes visibles. Estas composiciones presentan cuatro celdas con exceso en la primera, y ninguna en la segunda.

Una explicación plausible es que la heterogeneidad diluya la sensibilidad de un modelo concreto, porque agentes que responden de forma distinta a la misma señal se compensan al agregarse. Los datos no permiten confirmarla, aunque sí descartan la alternativa más inmediata. Los tres modelos que forman \textit{hetero\_local} presentan exceso al evaluarse por separado, de modo que la estabilidad del comité mixto no se explica porque sus miembros sean insensibles al atributo. Lo que permanece sin
resolver, y se detalla en la Sección~\ref{sec:amenazas_validez}, es si la diferencia procede de la diversidad o del reparto de modelos entre roles. Debe entenderse como una hipótesis que generan los resultados, no como una eficacia demostrada. Madigan et al.~\cite{madigan2025emergent} ya habían señalado que las propiedades de equidad de un sistema multiagente no se deducen de las de sus componentes; lo que aquí se añade es que tampoco se deducen manteniendo fijo el protocolo y cambiando solo la composición.

Ni el nivel de instrucción ni el atributo se comportan como cabría suponer. Enriquecer el rol no reduce el efecto de forma progresiva, y el nivel con identidad no actúa como estabilizador ya que en unas configuraciones el exceso permanece y en otras desaparece. Los dos atributos, por su parte, apuntan descriptivamente en direcciones opuestas, con 16 de 18 estimaciones negativas en geografía y 12 de 18 positivas en género. Ninguna de las dos regularidades supera el control de multiplicidad, de modo que no cabe hablar de penalización geográfica. Sí indican que ambas manipulaciones activan respuestas diferentes, algo esperable tratándose de intervenciones distintas.

El momento en que se mide resulta determinante. En génesis el agente ciego recibe exactamente la misma entrada en control y variante, lo que lo convierte en una referencia limpia. Desde el primer intercambio deja de serlo, porque empieza a recibir de forma indirecta las consecuencias del atributo a través de las respuestas de los agentes visibles. Por eso el ratio cae tras la primera ronda, aunque no porque la deliberación estabilice nada. Ambas dispersiones crecen, aumentando la visible un 4\,\% y la del agente ciego un 72\,\%. Al dispararse el denominador mucho más deprisa que
el numerador, el cociente disminuye pese a que las decisiones son en realidad menos consistentes que en génesis. El placebo confirma que buena parte de ese ensanchamiento pertenece a la deliberación misma y no al atributo. Que el fenómeno se concentre antes de cualquier comunicación coincide con la localización temporal que observan Nguyen et al.~\cite{nguyen2025social}. Sin embargo, en el presente estudio este patrón se observa sobre un desenlace continuo y con atributos ajenos a la identidad de los agentes. La transmisión posterior hacia el agente ciego existe como asociación, pero solo en \textit{homo\_mistral} supera lo que produce el placebo, y esa singularidad procede de una pendiente placebo anómalamente baja antes que de una capacidad de transmisión mayor. No hay base para hablar de un mecanismo general de propagación.

La mitigación por instrucciones apunta en la buena dirección sin llegar a demostrarse. Ocho de las doce celdas reducen la varianza visible y el patrón es consistente en geografía, pero solo una comparación excluye el valor de referencia con su intervalo de confianza. La campaña deja además dos advertencias de método. Las instrucciones alteran la variabilidad del propio agente ciego, de modo que una caída del ratio puede venir de la referencia y no de la decisión que se pretende estabilizar. Comparada con la heterogeneidad estructural, la intervención por \textit{prompt} muestra un efecto más débil, aunque ambas evidencias proceden de diseños distintos y no permiten declarar superior a ninguna de las dos.

En conjunto, lo observado se parece más a una sensibilidad inestable que a una preferencia. El sistema no favorece ni penaliza de forma sostenida a las empresas según el atributo, pero bajo ciertas configuraciones deja de tratarlas de manera consistente. Esa inestabilidad depende sobre todo de cómo esté compuesto el comité, no se corrige enriqueciendo las instrucciones y queda oculta a cualquier evaluación que se limite a comparar medias.

\section{Amenazas a la validez}
\label{sec:amenazas_validez}

La interpretación de los resultados depende, en primer lugar, de que los pares contrafactuales difieran únicamente en la manipulación estudiada. El diseño mantiene constantes los fundamentales, las empresas de contexto y la posición de la empresa sujeto, y sustituye nombres y \textit{tickers} por alias. La anonimización no es, aun así, completa. Una combinación suficientemente característica de sector, tamaño y perfil financiero podría permitir al modelo reconocer la compañía original, en cuyo caso asociaciones aprendidas sobre ella intervendrían en la decisión sin aparecer en la entrada.

Las dos manipulaciones no son además equivalentes en su significado. Cambiar el nombre de la persona que ocupa la dirección ejecutiva no guarda relación financiera con la ficha presentada. Cambiar la sede sí admite una lectura legítima, porque incluso con fundamentales idénticos una localización distinta puede evocar riesgos regulatorios, de tipo de cambio o institucionales que un inversor consideraría pertinentes. Los resultados geográficos se interpretan por ello como sensibilidad a la localización bajo información financiera controlada, y no como prueba de discriminación injustificada. El diseño aísla el efecto de presentar una sede distinta, pero no permite juzgar si las asociaciones que el modelo emplea para responder a ella son normativamente admisibles.

Una limitación análoga afecta al eje de niveles de instrucción. El nivel 2 reproduce literalmente la descripción funcional del nivel 1 y le añade una única oración de identidad por agente, de modo que la manipulación consiste en dotar al agente de una biografía y no en enriquecer su rol. Esas biografías sitúan a los tres agentes en centros financieros occidentales, y una de ellas declara además una especialización en narrativas de mercado occidentales. La comparación entre ambos niveles no separa por ello el efecto de asignar una identidad del de asignar esa procedencia concreta, distinción que resulta especialmente relevante para el eje geográfico. Deshacerla exigiría repetir el nivel con identidades equivalentes en extensión pero ancladas en otras regiones, según se plantea en la Sección~\ref{sec:trabajo_futuro}.

La referencia interna también presenta su propia amenaza. Que las entradas del agente ciego sean idénticas entre brazos, como se comprobó en el Capítulo~\ref{chap:validacion}, no implica que sus respuestas deban serlo, porque la generación sigue siendo estocástica y algunas configuraciones de API añaden fuentes adicionales de no determinismo. El riesgo se materializó en \textit{homo\_gpt} con nivel profesional, donde las dos estimaciones de variabilidad ciega resultaron incompatibles entre sí. Esas celdas se mantuvieron dentro de la familia de contrastes para no alterar retrospectivamente la corrección por multiplicidad, pero quedaron fuera de la interpretación. La validez del control cambia además con el tiempo, ya que desde $t=1$ el agente ciego observa a compañeros que sí recibieron la ficha completa. Los ratios posteriores a génesis describen por ello la dinámica del sistema y no repiten el contraste inicial.

El desenlace principal no fue preespecificado. El exceso de dispersión se adoptó tras analizar \textit{homo\_llama}, que conserva por ello un estatus exploratorio. Para acotar el riesgo de haber elegido la métrica que los datos favorecían, la definición del \textit{gap}, el ratio de varianzas y las familias de contrastes quedaron fijadas antes de examinar las cinco composiciones restantes, cuya evidencia constituye la parte prospectiva del análisis.

Los análisis estadísticos también dependen de supuestos que no pueden garantizarse por completo con este tamaño muestral. La comparación principal de varianzas utiliza el contraste de Fisher, sensible a desviaciones de la normalidad, por lo que se complementó con un contraste robusto basado en la mediana. Diez de las doce celdas con exceso de dispersión interpretable reciben apoyo de ambos procedimientos, mientras que dos dependen únicamente del contraste principal. Los análisis del \textit{gap} medio presentan otras restricciones: la prueba por inversión de signos supone una distribución aproximadamente simétrica de los \textit{gaps}, y los intervalos \textit{bootstrap} se estiman a partir de un número limitado de pares. La corrección de Benjamini--Hochberg reduce el riesgo de falsos positivos asociado a realizar múltiples contrastes, pero no aumenta la capacidad para detectar efectos pequeños. Por ello, la interpretación se apoya en la consistencia entre varios análisis y no en un único valor de $p$.

Dos elecciones de diseño condicionan finalmente lo que puede afirmarse. La primera afecta a la interpretación del contraste entre composiciones. Los tres modelos que forman \textit{hetero\_local} presentan exceso de dispersión cuando se evalúan por separado en comités homogéneos, por lo que la mayor estabilidad del comité mixto no puede explicarse simplemente porque sus miembros sean insensibles al atributo. Sin embargo, en las composiciones homogéneas cada modelo ocupa los tres roles, mientras que en \textit{hetero\_local} cada uno desempeña únicamente uno. Que un modelo muestre sensibilidad cuando ocupa todos los puestos no implica que la presente también en el rol concreto que asume dentro del comité heterogéneo. La comparación con \textit{hetero\_api\_gpt} es menos directa, ya que \textit{homo\_gpt} presenta problemas de estabilidad en la referencia ciega
que limitan la interpretación de parte de sus resultados. El contraste entre \textit{hetero\_local} y \textit{hetero\_api\_gpt} permite estudiar el efecto de sustituir únicamente al segundo agente,
manteniendo fijos los otros dos. Sin embargo, atribuir la mayor estabilidad de los comités heterogéneos a la diversidad como propiedad general requeriría evaluar más combinaciones y permutaciones de modelos entre los distintos roles. La segunda elección afecta a cómo se define la decisión colectiva. La asignación del comité se calcula como la media de los dos agentes visibles y no procede de una regla explícita de consenso, votación o negociación. Las conclusiones sobre estabilidad colectiva dependen, por tanto, de esta forma de agregar las respuestas, ya que un mecanismo distinto podría reducir o aumentar las discrepancias entre los agentes aun partiendo de las mismas decisiones individuales.

\section{Limitaciones}
\label{sec:limitaciones}

Los resultados deben interpretarse dentro del alcance del banco de pruebas construido. La campaña emplea cuatro modelos de lenguaje, tres de ellos abiertos y ejecutados localmente y uno accesible mediante API, combinados en seis composiciones. La matriz prevista contemplaba ocho, pero las dos que incorporaban un modelo comercial adicional quedaron fuera por restricciones de presupuesto, según se detalla en la Sección~\ref{sec:presupuesto}. La selección permite contrastar arquitecturas y procedencias distintas, aunque no representa la variedad de modelos disponibles ni autoriza a suponer que el patrón se reproduzca en modelos mayores, en versiones posteriores o bajo otros procedimientos de alineamiento. Todos los comités constan además de tres agentes con una estructura fija de roles; sistemas con más participantes, otras topologías de comunicación o reglas explícitas de coordinación podrían comportarse de otro modo.

El escenario se limita a una única clase de decisión financiera. Los agentes reparten un presupuesto fijo entre empresas a partir de fundamentales, información contextual y noticias, sin que la calidad de la asignación se evalúe mediante rentabilidad futura ni mediante su desempeño en un mercado real. La tarea funciona como entorno controlado para medir sensibilidad contrafactual y no como simulación del proceso de inversión. De los resultados no se sigue, por tanto, que el exceso de dispersión produzca pérdidas económicas, ni que reaparezca en decisiones financieras de otra naturaleza como la concesión de crédito o la fijación de precios.

El conjunto de intervenciones contrafactuales es igualmente reducido, y su construcción incorpora simplificaciones deliberadas. Se manipulan dos atributos de la empresa, el género señalado mediante la identidad de la dirección ejecutiva y la localización de la sede, sobre un conjunto acotado de perfiles y sectores. El género se opera de forma binaria y las sedes no occidentales se agrupan en una única categoría, decisiones tomadas para conservar potencia estadística que reducen la granularidad de lo que puede afirmarse y cuyas implicaciones se retoman en la Sección~\ref{sec:consideraciones_eticas}. Los resultados geográficos dependen además de las localizaciones concretas empleadas y no se extienden a cualquier diferencia entre países o regiones.

La profundidad experimental también acota las conclusiones. Cada configuración se evalúa con tres semillas sobre veintiuna parejas contrafactuales, y la deliberación se detiene tras cuatro rondas. Un número mayor de repeticiones permitiría estimar con más precisión las colas de las distribuciones y detectar efectos pequeños, y deliberaciones más largas podrían revelar patrones de convergencia o amplificación ajenos a este horizonte. A ello se añade que la referencia placebo se ejecutó una vez por composición, con el nivel profesional, y no para cada celda de la matriz. No existe por tanto evidencia de que ese suelo de ruido externo se mantenga invariable entre niveles de instrucción, motivo por el cual los contrastes principales se apoyan en la referencia ciega correspondiente a cada celda.

La campaña de mitigación tiene el alcance más estrecho de todos. Las dos instrucciones se evalúan únicamente sobre \textit{homo\_qwen}, en los tres niveles y en el turno de génesis. Esa elección permite comprobar si una intervención ligera modifica un sistema que había mostrado exceso de forma recurrente, pero no generalizar su eficacia a otros modelos. La restricción al turno inicial, motivada por la localización del efecto en la decisión de génesis, impide además evaluar el componente que distingue a la instrucción activa, cuyo mecanismo solo puede operar cuando el agente observa el razonamiento del resto del comité. La comparación entre ambas queda limitada a comprobar la sensibilidad del efecto a la redacción del \textit{prompt}; evaluar la corrección activa entre agentes exigiría repetir la campaña con deliberación completa, según se plantea en la Sección~\ref{sec:trabajo_futuro}. Tampoco se estudian intervenciones que requieran modificar los parámetros del modelo ni mecanismos externos de supervisión.

Estas limitaciones delimitan el objetivo del estudio. El banco de pruebas no aspira a ofrecer una estimación universal del sesgo de los modelos de lenguaje ni una estrategia de mitigación válida para cualquier sistema multiagente. Su aportación consiste en caracterizar, bajo condiciones controladas y reproducibles, una forma de sensibilidad que puede pasar inadvertida cuando solo se examinan diferencias medias, y en señalar qué dimensiones del fenómeno requieren validación en modelos, dominios y arquitecturas más amplios.